

\subsubsection*{Name the tax liability settled by each payment}
Once carried tax and newly assessed tax are both payable, the payment amount
alone cannot identify which liability has been settled. A payment of 3 can reduce
either the carried payable of $19/4$ to $7/4$, leaving the new payable of 5
unchanged, or reduce the new payable to 2, leaving the carried payable at $19/4$.
Both choices reduce cash and total payable by 3. Their remaining tax obligations
have different histories, which the simulator must preserve.

The continuation therefore requires the caller to name each tax vintage and its
allocated amount. Every amount must be positive, no greater than the named
outstanding balance, and the allocations must sum to the payment. Actual shared
cash must cover that sum. An unknown authorization, a duplicated vintage, an
unknown liability or an excessive allocation prevents that operation. The
method supplies no first-in-first-out convention or proportional split.

For example, take the hypothetical first assessment through 30 June 2026 inside
the financial period ending on 31 December. Carried tax of $19/4$ is still
unpaid, new tax is 5, cash is $721/2$, and recorded profit is $-15$. Allocate 1
of a payment to carried tax and 2 to new tax. Cash becomes $715/2$; the two
payables become $15/4$ and 3, totaling $27/4$. Tax expense remains 5 and profit
remains $-15$. The payment is recorded in both affected settlement histories,
with their distinct amounts and the same dated cash movement.

An ordinary event can follow only from 1 July, the assessment's exclusive
endpoint. A further right-of-use depreciation charge of 2 reduces its asset from
80 to 78 and profit to $-17$, without changing either tax payable. Assets then
total $1367/2$, liabilities $467/4$, and equity including profit $2267/4$.
Operating cash flow remains $-3$. This depreciation is explicitly listed as
activity after the assessment; it has not received a new tax assessment.

The full book, prior loss inventory, deferred balances and unresolved findings
remain attached to the result. If a later operation fails, earlier completed
events and their settlement histories survive as a recorded prefix. Payments
can occur on the recorded assessment day, but an assessment at the financial
period's endpoint leaves no later date for ordinary events in that period.
Cumulative reassessment, new-period deferred recognition and another period
carry require their own compatible methods. The present calculation supplies
conditional bookkeeping evidence rather than annual accounting acceptance.
